% Problem: limits of quotients (English).

\begin{exercise}[Limits of quotients]
Compute the following limits:
\[
  \text{(i)}\ \lim_{x \to 2} \frac{x^2 - 4}{x - 2} ,
  \qquad
  \text{(ii)}\ \lim_{x \to 1} \frac{x^3 - 1}{x^2 - 1} ,
  \qquad
  \text{(iii)}\ \lim_{x \to 0} \frac{\sqrt{1 + x} - 1}{x} .
\]

\dmarks{3}

\begin{hint}
In all three, numerator and denominator vanish at the point. Factor to cancel
the common factor; in the last one, multiply and divide by
$\sqrt{1 + x} + 1$.
\end{hint}

\begin{answer}
(i) $4$; (ii) $3/2$; (iii) $1/2$.
\end{answer}

\begin{solution}
Since the limit does not depend on the value at the point, we may simplify
for $x$ different from it and evaluate afterwards.
\begin{parts}
  \item For $x \neq 2$, $\dfrac{x^2 - 4}{x - 2} = x + 2$, so the limit is
        $4$.
  \item For $x \neq \pm 1$,
        \[
          \frac{x^3 - 1}{x^2 - 1}
          = \frac{(x - 1)(x^2 + x + 1)}{(x - 1)(x + 1)}
          = \frac{x^2 + x + 1}{x + 1} ,
        \]
        which tends to $3/2$.
  \item For $x \neq 0$,
        \[
          \frac{\sqrt{1 + x} - 1}{x}
          = \frac{(1 + x) - 1}{x\,\bigl(\sqrt{1 + x} + 1\bigr)}
          = \frac{1}{\sqrt{1 + x} + 1} ,
        \]
        which tends to $1/2$.
\end{parts}
\end{solution}

\begin{marking}
One mark per part. Writing ``$0/0$'' as the result scores nothing; a correct
simplification with a slip in the final evaluation, half a mark.
\end{marking}
\end{exercise}
